\subsection{Estado industrial y normativo: qué existe ya en producto}
\label{subsec:tf-industrial}

Una brecha académica no basta para justificar un trabajo de máster en robótica
aplicada: si la industria ya resolviera el problema con producto documentado, la
contribución sería tardía aunque la literatura arbitrada la ignorase. La revisión
incorpora por eso catorce precedentes comerciales seleccionados con el muestreo
intencional de la Sección~\ref{subsec:method-review-protocol}. La
Figura~\ref{fig:tf-industrial-map} los sitúa sobre dos dimensiones ordinales:
interacción cooperativa entre vehículos y descentralización de la coordinación.

% Figura 1 del consolidado industrial, portada literalmente desde
% `thesis-reviews/TFM_estado_industrial_patentes_EDITABLE_v47.tex` (líneas
% 148-247). Solo se cambia el envoltorio: el original usa \figcap y aquí va
% dentro de un `figure` con el pie y la fuente en el estilo de la memoria.
% El cuerpo TikZ y la paleta son los del autor, de modo que sigue siendo
% editable con la misma sintaxis que en el consolidado.
\begin{figure}[!htbp]
  \centering
  \makebox[\textwidth][c]{\resizebox{\textwidth}{!}{%
\begin{tikzpicture}[x=1mm,y=1mm,font=\sffamily\fontsize{6.45}{6.95}\selectfont]
% La caja original (0,-14)--(174,98) dejaba fuera las etiquetas ancladas al este
% del borde derecho, que al reescalar invadían el margen. Se amplía para que
% \resizebox mida el contenido completo.
\path[use as bounding box] (0,-14) rectangle (196,98);

% plot frame and regions: four clear quadrants
\fill[gray!8] (30,16) rectangle (98,53);
\fill[orange!10] (98,16) rectangle (166,53);
\fill[teal!10] (30,53) rectangle (98,90);
\fill[navy!8] (98,53) rectangle (166,90);
\draw[rule,line width=.45pt] (30,16) rectangle (166,90);
\draw[rule!88,line width=.42pt] (98,16)--(98,90);
\draw[rule!88,line width=.42pt] (30,53)--(166,53);
\foreach \x in {64,132}{\draw[rule!75,densely dotted,line width=.24pt] (\x,16)--(\x,90);}
\foreach \y in {34.5,71.5}{\draw[rule!75,densely dotted,line width=.24pt] (30,\y)--(166,\y);}

% axes
\node[anchor=north,align=center,text width=30mm] at (47,14.7){Flota\ independiente};
\node[anchor=north,align=center,text width=30mm] at (81,14.7){Movimiento\\ coordinado};
\node[anchor=north,align=center,text width=30mm] at (115,14.7){Misma carga\\ sincronizada};
\node[anchor=north,align=center,text width=30mm] at (149,14.7){Acoplamiento\\ rígido/modular};
\node[anchor=north,font=\bfseries] at (98,5.2){Interacción cooperativa entre vehículos};
\node[anchor=north,font=\sffamily\fontsize{5.5}{5.8}\selectfont,text=muted] at (98,9.4){menor $\rightarrow$ mayor};

\node[anchor=east,align=right,text width=26mm] at (27,25){Humano /\\ configurado};
\node[anchor=east,align=right,text width=26mm] at (27,43.5){Control\\ central};
\node[anchor=east,align=right,text width=26mm] at (27,62){Líder /\\ master};
\node[anchor=east,align=right,text width=26mm] at (27,80.5){Entre vehículos /\\ distribuido};
\node[rotate=90,anchor=south,font=\bfseries] at (4.5,53){Descentralización de la coordinación y el control};
\node[rotate=90,anchor=south,font=\sffamily\fontsize{5.3}{5.6}\selectfont,text=muted] at (9.6,53){menor $\rightarrow$ mayor};

% label and point styles
\tikzset{regtitle/.style={anchor=west,font=\bfseries,fill=white,fill opacity=.84,text opacity=1,inner sep=1.1pt,rounded corners=.8pt}, regsub/.style={anchor=west,text=muted,fill=white,fill opacity=.84,text opacity=1,inner sep=.9pt,rounded corners=.8pt}, pt/.style={circle,draw=onyx,line width=.55pt,minimum size=5.2mm,inner sep=0pt,fill=white}, num/.style={font=\sffamily\fontsize{5.8}{6.0}\selectfont\bfseries,text=white}, ctx/.style={diamond,aspect=1.15,draw=onyx!65,line width=.45pt,minimum width=5.0mm,minimum height=4.3mm,inner sep=0pt,fill=onyx!14!white}, ctxnum/.style={font=\sffamily\fontsize{5.1}{5.3}\selectfont\bfseries,text=onyx}}

% region labels
\node[regtitle,text=teal!90!black] at (34,87.8){flota autónoma};
\node[regsub] at (34,84.4){coordinación distribuida; sin carga compartida};
\node[regtitle,text=navy] at (101,87.8){convergencia objetivo};
\node[regsub] at (101,84.4){carga compartida + mayor autonomía};
\node[regtitle,text=ink] at (34,20.8){automatización convencional};
\node[regsub] at (34,17.6){1 carga por robot; control configurado o central};
\node[regtitle,text=orange!90!black] at (100.5,18.9){heavy-load configurado};
\node[regsub] at (100.5,15.9){carga compartida/acople; baja autonomía};

% points + labels
\node[pt,fill=orange] (c1) at (151,25) {}; \node[num] at (c1){1};
\draw[orange!70,line width=.32pt] (c1)--(141.5,23.3); \node[anchor=east,align=right] at (141.0,23.3){KUKA omniMove / Airbus\\2016};

\node[pt,fill=teal] (c2) at (136.5,79.0) {}; \node[num] at (c2){2};
\draw[teal!70,line width=.32pt] (c2)--(128.8,75.4); \node[anchor=east,align=right] at (128.3,75.4){Stäubli + R3\\2022};

\node[pt,fill=blue] (c3) at (116,47) {}; \node[num,text=onyx] at (c3){3};
\draw[blue!70,line width=.32pt] (c3)--(109.5,55.5); \node[anchor=east,align=right] at (109.0,55.5){SIASUN dual-vehicle\\2024};

\node[pt,fill=spicy!68!white] (c4) at (113,63) {}; \node[num] at (c4){4};
\draw[plum!70,line width=.32pt] (c4)--(103,67); \node[anchor=east,align=right] at (102.5,67){KUKA KMP 3000P\\2025};

\node[pt,fill=orange] (c5) at (136,29.5) {}; \node[num] at (c5){5};
\draw[orange!70,line width=.32pt] (c5)--(127.5,28.8); \node[anchor=east,align=right] at (127.0,28.8){GREATOX GTVL25\\cat. 2026};

\node[pt,fill=orange] (c6) at (158,32) {}; \node[num] at (c6){6};
\draw[orange!70,line width=.32pt] (c6)--(165,39.0); \node[anchor=west,align=left] at (165.5,39.0){TII SCHEUERLE SPMT\\2023};

\node[pt,fill=teal] (c7) at (48,83) {}; \node[num] at (c7){7};
\draw[teal!70,line width=.32pt] (c7)--(56,78.5); \node[anchor=west,align=left] at (56.5,78.5){AGILOX X-SWARM / BMW\\2023};

\node[pt,fill=spicy!68!white] (c8) at (131,64) {}; \node[num] at (c8){8};
\draw[spicy!70,line width=.32pt] (c8)--(140.0,59.0); \node[anchor=west,align=left] at (140.5,59.0){Beacon Dual-AMR\\2026};

\node[pt,fill=orange] (c9) at (148,38) {}; \node[num] at (c9){9};
\draw[orange!70,line width=.32pt] (c9)--(137.5,47.0); \node[anchor=east,align=right] at (137.0,47.0){HUBTEX Aviation\\2020};

\node[pt,fill=blue] (c10) at (53,51.5) {}; \node[num,text=onyx] at (c10){10};
\draw[blue!70,line width=.32pt] (c10)--(60,54); \node[anchor=west,align=left] at (60.5,54){Hyundai HMGICS\\2023};

% contextual central-fleet systems (not expanded in Table 1)
\node[ctx] (ctxA) at (40.5,42.0) {}; \node[ctxnum] at (ctxA){A};
\draw[mustard!80!onyx,line width=.28pt] (ctxA)--(35.0,38.5); \node[anchor=east,align=right,font=\sffamily\fontsize{5.7}{6.0}\selectfont] at (34.5,38.5){MiR Fleet};

\node[ctx] (ctxB) at (58.0,39.0) {}; \node[ctxnum] at (ctxB){B};
\draw[mustard!80!onyx,line width=.28pt] (ctxB)--(61.5,35.0); \node[anchor=west,align=left,font=\sffamily\fontsize{5.7}{6.0}\selectfont] at (62.0,35.0){OTTO Fleet\\Manager};

\node[ctx] (ctxC) at (74.0,46.0) {}; \node[ctxnum] at (ctxC){C};
\draw[mustard!80!onyx,line width=.28pt] (ctxC)--(77.0,41.8); \node[anchor=north,align=center,font=\sffamily\fontsize{5.7}{6.0}\selectfont] at (77.0,41.4){Geek+ RMS};

\node[ctx] (ctxD) at (89.0,40.5) {}; \node[ctxnum] at (ctxD){D};
\draw[mustard!80!onyx,line width=.28pt] (ctxD)--(92.0,36.5); \node[anchor=west,align=left,font=\sffamily\fontsize{5.7}{6.0}\selectfont] at (92.5,36.5){Swisslog\\IntraMove};

% convergence vectors and TFM target
% The target is not placed at the rigid-coupling extreme: the TFM seeks carga compartida
% multi-contact cooperation with distributed coalition decisions, not necessarily a rigid robot-robot module.
\node[rounded corners=1.0pt,minimum size=6.4mm,inner sep=0pt,draw=spicy,line width=.55pt,fill=autumn] (tfm) at (147.5,84.0) {};
\node[anchor=west,font=\bfseries,text=spicy] at (152.2,84.0){TFM objetivo};

% compact legend
\node[anchor=north west,align=left,text width=162mm,text=muted] at (5,1.5){\fontsize{6.4}{7}\selectfont Forma: círculo = caso principal; rombo gris = contexto de flota; cuadrado naranja = TFM.\\ \textbf{En los círculos}, el color indica: \textcolor{orange}{\textbf{configurado/humano}}, \textcolor{blue}{\textbf{central}}, \textcolor{plum}{\textbf{líder/seguidor}}, \textcolor{teal}{\textbf{entre vehículos/distribuido}}.};
\end{tikzpicture}%
  }}
  \caption[Precedentes industriales de cooperación multi-vehículo.]{Diez
  precedentes principales y cuatro referencias contextuales de gestión de flota,
  clasificados según dos dimensiones ordinales: interacción cooperativa entre
  vehículos (izquierda $\rightarrow$ derecha) y descentralización de la
  coordinación y el control (abajo $\rightarrow$ arriba). Los círculos numerados
  corresponden a la Tabla~\ref{tab:tf-industrial-matrix}; los rombos A--D son
  contexto de flota centralizada. La posición dentro de una categoría se desplaza
  solo para evitar solapes y no tiene interpretación métrica.}
  \label{fig:tf-industrial-map}
  \viusource{elaboración propia a partir de las fuentes oficiales citadas en la
  Tabla~\ref{tab:tf-industrial-matrix}}
\end{figure}


La lectura relevante es que dos líneas funcionales coexisten con precedentes
comerciales propios. En el cuadrante inferior derecho, la cooperación física para carga
pesada está resuelta con configuración de ingeniería y escasa autonomía: KUKA omniMove
mueve secciones de fuselaje de hasta 100~toneladas en Airbus Stade con composición
dedicada \parencite{kukaOmnimoveAirbus2016} y los transportadores modulares SPMT
resuelven \emph{project cargo} mediante configuración previa
\parencite{tiiScheuerleSpmt2023}. En el superior izquierdo, la autonomía de flota está
igualmente resuelta pero sin carga compartida: AGILOX opera 27~AMR en BMW Regensburg con
coordinación entre pares y sin control central \parencite{agiloxXswarmBmw2023}. El
cuadrante superior derecho, donde ambas capacidades convergen, permanece escasamente
poblado.

% Tabla 1 del consolidado industrial, portada desde
% `thesis-reviews/TFM_estado_industrial_patentes_EDITABLE_v47.tex` (líneas
% 263-334). Cambios respecto del original, todos de integración:
%   - los tipos de columna L{}/C{}/Y pasan a J{}/K{}/Z, porque C e Y ya
%     existen en viu-mrob-thesis.sty con otra signatura;
%   - el hueco de fotografía se sustituye por el número de caso, que es lo
%     que casa la fila con la Figura del mapa; el enlace a la fuente oficial
%     se conserva sobre el nombre del sistema;
%   - las marcas \casecite se sustituyen por citas biblatex reales.
\begin{table}[!htbp]
  \centering
  \caption{Matriz funcional de precedentes industriales frente a la cadena operativa del trabajo.}
  \label{tab:tf-industrial-matrix}
{\setstretch{1}\fontsize{7.35}{8.10}\selectfont
 \setlength{\tabcolsep}{0.4pt}
 \renewcommand{\arraystretch}{1.32}
\begin{tabularx}{\linewidth}{@{}J{3.72cm}K{0.74cm}K{0.78cm}K{0.86cm}K{0.86cm}K{0.82cm}K{0.90cm}Z@{}}
\toprule
\textbf{Sistema / evidencia} &
{\fontsize{4.9}{5.2}\selectfont\shortstack{\textbf{Carga}\\\textbf{comp.}}} &
{\fontsize{4.9}{5.2}\selectfont\shortstack{\textbf{Multi-veh.}\\\textbf{($>2$)}}} &
{\fontsize{4.9}{5.2}\selectfont\shortstack{\textbf{Selecc.}\\\textbf{miembros}}} &
{\fontsize{4.9}{5.2}\selectfont\shortstack{\textbf{Coord./ctrl.}\\\textbf{local}}} &
{\fontsize{4.9}{5.2}\selectfont\shortstack{\textbf{Cert.}\\\textbf{factib.}}} &
{\fontsize{4.9}{5.2}\selectfont\shortstack{\textbf{Recup.}\\\textbf{coalición}}} &
\textbf{Evidencia industrial / límite}\\
\midrule
\rowcolor{soft}
\matrixsystem{1}{KUKA · omniMove}{Airbus Stade}{2016}{https://www.kuka.com/en-us/industries/solutions-database/2016/07/solution-robotics-airbus}
& \capyes & \capunk & \capno & \capno & \capno & \capno
& 100 t; composición dedicada. Sin selección dinámica, certificado online ni recomposición autónoma publicada \parencite{kukaOmnimoveAirbus2016}.\\

\matrixsystem{2}{Stäubli + R3}{Virtual drawbar}{2022}{https://www.r3.group/en/success-stories/collaborative-agvs-in-the-cleanroom}
& \capyes & \capno & \capno & \cappart & \capno & \capno
& 2 AGV, 2$\times$5 t; ejecución V2V directa a 32 ms con PROFINET/PROFIsafe. La fuente no demuestra formación distribuida de la pareja ni ausencia de una autoridad dominante \parencite{staubliR3Drawbar2022}.\\

\rowcolor{soft}
\matrixsystem{3}{SIASUN}{Heavy-truck dual vehicle}{2024}{https://en.siasun.com/heavy-truck-assembly-line-mobile-robot.html}
& \capyes & \capno & \cappart & \capno & \capno & \capno
& Dual-vehicle linkage + MES. Selección a nivel de planta, no reclutamiento local ni recomposición física \parencite{siasunDualVehicle2024}.\\

\matrixsystem{4}{KUKA · KMP 3000P}{Load sharing}{2025}{https://www.kuka.com/en-us/products/amr-autonomous-mobile-robotics/mobile-platforms/kmp-3000p-omnimove}
& \capyes & \capno & \capno & \cappart & \capno & \capno
& Dos plataformas, hasta 6 t; líder--seguidor. Carga compartida comercial, pero pareja fija \parencite{kukaKmp3000p2025}.\\

\rowcolor{soft}
\matrixsystem{5}{GREATOX · GTVL25}{Multi-vehicle synchronous}{cat. 2026}{https://en.greatoxagv.com/products_info/id-14.html}
& \capyes & \cappart & \capno & \capunk & \capno & \capunk
& 80--220 t y cooperación multi-vehículo. Monitoriza centro de gravedad; no publica certificado online de factibilidad mecánica, selección o recuperación \parencite{greatoxGtvl2026}.\\

\matrixsystem{6}{TII SCHEUERLE · SPMT}{Modular heavy transport}{2023}{https://www.tii-group.com/tii-scheuerle/our-solutions/spmt/scheuerle-spmt}
& \capyes & \cappart & \capno & \capno & \capno & \capno
& \textit{Project cargo} modular; configuración de ingeniería, no certificado online ni reclutamiento autónomo \parencite{tiiScheuerleSpmt2023}.\\

\rowcolor{soft}
\matrixsystem{7}{AGILOX · X-SWARM}{BMW Regensburg}{2023}{https://www.agilox.net/en/blog/bmw-agilox-swarm/}
& \capno & \capna & \capna & \capyes & \capna & \capna
& 27 AMR; coordinación entre pares sin líder central. Una carga por AMR: no aplica coalición física \parencite{agiloxXswarmBmw2023}.\\

\matrixsystem{8}{Beacon Robot · Dual-AMR}{Collaborative transport}{2026}{https://www.beacon-robot.com/news/beacon-robot-amr-dual-robot-collaboration-officially-released-from-working-alone-to-working-together/}
& \capyes & \capno & \capyes & \cappart & \capno & \capno
& Emparejamiento automático y relevo de autoridad ante fallo documentados. 2 robots; despacho central y líder--seguidor. No se documenta sustitución/reclutamiento de un nuevo miembro \parencite{beaconDualAmr2026}.\\

\rowcolor{soft}
\matrixsystem{9}{HUBTEX · SL/HL-AGV}{Aviation}{2020}{https://www.hubtex.com/sites/default/files/2020-04/HUBTEX_Aircraft\%20Components_ENG.pdf}
& \capyes & \capunk & \capno & \capno & \capno & \capno
& Hasta 70 t; plataformas acoplables y operación semi/remota. Sin selección dinámica, certificado online o recomposición \parencite{hubtexAviation2020}.\\

\matrixsystem{10}{Hyundai · HMGICS}{AGV / AMR}{2023}{https://www.hyundaimotorgroup.com/en/story/CONT0000000000126181}
& \capno & \capna & \capna & \cappart & \capna & \capna
& Orquestación de planta y reactividad local; sin coalición física de carga compartida \parencite{hyundaiHmgics2023}.\\
\midrule
\rowcolor{softplum}
\multicolumn{1}{@{}l}{\parbox[t]{3.55cm}{\textbf{TFM objetivo}\\[-0.3pt]{\fontsize{5.35}{5.65}\selectfont\color{muted}(referencia conceptual)}}}
& \capgoal & \capgoal & \capgoal & \capgoal & \capgoal & \capgoal
& Integrar carga compartida, coalición heterogénea de tamaño adaptable, reclutamiento local, coordinación/control distribuido, verificación de factibilidad contacto/torsor y recomposición tras fallo.\\
\bottomrule
\end{tabularx}}

\vspace{2.6pt}
\noindent\colorbox{softwarm}{\parbox{0.972\linewidth}{\setstretch{1}\fontsize{7.35}{8.1}\selectfont
\textbf{Lectura de la matriz.} La industria demuestra por separado carga compartida, emparejamiento/relevo ante fallo y coordinación distribuida de flota. \textbf{Ninguno de los precedentes revisados documenta simultáneamente una coalición física heterogénea multi-vehículo, selección local de miembros, verificación mecánica de factibilidad y recomposición durante la misión.}}}

\vspace{1.8pt}
{\setstretch{1}\fontsize{7.15}{8.0}\selectfont\color{muted}
Contexto A--D del mapa: MiR Fleet, OTTO Fleet Manager, Geek+ RMS y Swisslog IntraMove coordinan flotas de carga individual; no se incluyen como filas principales porque no aportan carga compartida. Los símbolos indican evidencia pública revisada, no una puntuación de similitud. En Multi-veh. ($>2$), \cappart\ indica cooperación multi-vehículo documentada sin cardinalidad superior a dos explícita en la fuente.\par}

\vspace{2pt}
{\centering\footnotesize
\capyes\ documentado \quad \cappart\ parcial o condicionado \quad
\capno\ no documentado públicamente \quad \capunk\ arquitectura no publicada \quad
\capna\ no aplica \quad \capgoal\ objetivo del TFM.\\
«No documentado» no equivale a inexistencia.\par}
  \viusource{elaboración propia con documentación técnica oficial de cada fabricante}
\end{table}


La Tabla~\ref{tab:tf-industrial-matrix} convierte el mapa en una comparación capacidad a
capacidad contra la cadena operativa de este trabajo. Beacon Dual-AMR es el precedente
con mayor solapamiento funcional, porque documenta emparejamiento automático y relevo de
autoridad ante fallo, pero opera con dos robots bajo despacho central y no sustituye a
un miembro por otro de la flota \parencite{beaconDualAmr2026}. Conviene además separar
dos heterogeneidades: la de flota, entendida como convivencia de modelos distintos bajo
un mismo gestor, aparece con frecuencia en las plataformas de contexto; la que se da
dentro de una misma coalición física, donde robots con capacidades y geometrías
diferentes deben cerrar conjuntamente restricciones de carga, contacto y \emph{wrench},
tiene cobertura pública mucho menor y es la que el problema focal aborda.

% Tabla 2 del consolidado industrial, portada desde
% `thesis-reviews/TFM_estado_industrial_patentes_EDITABLE_v47.tex` (líneas
% 498-509). Las marcas \casecite se sustituyen por citas biblatex reales y
% los tipos de columna L{}/Y pasan a J{}/Z.
\begin{table}[!htbp]
  \centering
  \caption{Normas e interfaces relevantes para la integración industrial.}
  \label{tab:tf-standards}
{\setstretch{1}\fontsize{6.95}{7.65}\selectfont\setlength{\tabcolsep}{2.2pt}\renewcommand{\arraystretch}{1.16}
\begin{tabularx}{\linewidth}{@{}J{3.15cm}Z@{}}
\toprule
\textbf{Documento} & \textbf{Relación con el problema}\\
\midrule
\rowcolor{soft} VDA 5050 v3.0.0 (2026) & Interfaz de órdenes y estado entre control de flota y robots móviles. La v3.0 incorpora zonas de libre navegación y compartición de rutas (\textit{path sharing}); favorece interoperabilidad operativa, pero no define la constitución ni la verificación mecánica de una coalición de carga compartida \parencite{vda5050v3}.\\
ISO 21423 (2026, en publicación) & Protocolos de comunicación para interoperabilidad entre AMR multivendor, gestores de flota y recursos empresariales. Consultado el 11-09-2026: figura en etapa 60.00, ``International Standard under publication'', y excluye requisitos de seguridad \parencite{iso21423}.\\
\rowcolor{soft} ISO 3691-4:2023 & Requisitos de seguridad y medios de verificación para vehículos industriales sin conductor y sus sistemas; incluye AGV y AMR como ejemplos \parencite{iso36914}.\\
ANSI/A3 R15.08-3:2026 & Uso seguro de aplicaciones IMR: evaluación de riesgos, gestión de cambios y mantenimiento del riesgo aceptable durante el ciclo de vida \parencite{ansiA3R1508part3}.\\
\bottomrule
\end{tabularx}}
  \viusource{elaboración propia a partir de los documentos citados}
\end{table}


La Tabla~\ref{tab:tf-standards} delimita el marco normativo de cualquier realización
industrial de esta arquitectura. Ni VDA 5050 ni ISO 21423 definen la constitución de una
coalición física o su certificación mecánica, e ISO 3691-4 y ANSI/A3 R15.08-3 fijan el
marco de seguridad sin cubrir ese contrato. Queda por tanto fuera de la normativa
vigente y debe resolverlo el sistema, con una consecuencia de diseño directa: el
certificado de \emph{wrench} no puede delegarse en una capa de interoperabilidad y debe
producirse dentro del propio mecanismo de reclutamiento, como se formula en el
Capítulo~\ref{sec:methodology}.
